\setchapterpreamble[u]{\margintoc}
\chapter{Gradient-Based Optimisation}
\labch{gradient-methods}

\sources{Hass, Heil and Weir, \emph{Thomas' Calculus}, 14th ed.\ \cite{thomas2018calculus},
Chapter~14; Strang, \emph{Introduction to Linear Algebra}, 6th ed.\
\cite{strang2023ila}, Chapter~9; Zhang, Lipton, Li and Smola,
\emph{Dive into Deep Learning} \cite{d2l2023}, Chapters~11--12.}

This chapter closes Part~II by combining the two halves already assembled: the
gradient says which way to move, and the Hessian says how far the move can be
trusted.

\section{Descent directions and the basic method}
\labsec{descent}

By \refsec{steepest}, \(-\nabla f\) is the direction along which \(f\) falls
fastest, so the simplest possible method is
\begin{equation}
	\vect{x}_{k+1}=\vect{x}_k-\alpha\,\nabla f(\vect{x}_k),
	\labeq{gd}
\end{equation}
with step size \(\alpha>0\). Everything interesting is in the choice of
\(\alpha\), and \refeq{taylor2} explains why: the linear model that justified the
direction is only accurate while the quadratic term stays small.

\section{The quadratic case, exactly}
\labsec{gd-quadratic}

Take \(f(\vect{x})=\tfrac12\vect{x}\tp A\vect{x}-\vect{b}\tp\vect{x}\) with \(A\)
symmetric positive definite, whose minimiser \(\vect{x}^{\star}\) solves
\(A\vect{x}=\vect{b}\) by \refsec{quadratic-forms}. Since
\(\nabla f(\vect{x})=A\vect{x}-\vect{b}=A(\vect{x}-\vect{x}^{\star})\), the error
\(\vect{e}_k=\vect{x}_k-\vect{x}^{\star}\) obeys
\[
	\vect{e}_{k+1}=(I-\alpha A)\,\vect{e}_k .
\]

Diagonalising \(A=Q\Lambda Q\tp\) decouples the coordinates: along the
eigenvector \(\vect{q}_i\) the error is multiplied by \(1-\alpha\lambda_i\) at
every step.

\begin{theorem}
Gradient descent on this quadratic converges for every starting point if and
only if \(0<\alpha<2/\lambda_{\max}\). The best rate is obtained at
\(\alpha^{\star}=2/(\lambda_{\min}+\lambda_{\max})\), where the error contracts
by a factor
\[
	\frac{\kappa-1}{\kappa+1},
	\qquad
	\kappa=\frac{\lambda_{\max}}{\lambda_{\min}},
\]
per step.
\end{theorem}

\marginnote{With \(\kappa=100\) the factor is \(0.980\); reducing the error by
\(10^{-3}\) then takes about \(340\) steps. Conditioning, introduced in
\refsec{conditioning} as a numerical concern, is here a direct statement about
training time.}

\begin{remark}
The geometry is the elongated ellipse of \refsec{quadratic-forms}. The negative
gradient points across the valley rather than along it, so the iterates zigzag.
Nothing about this depends on the objective being artificial --- an
ill-conditioned Hessian produces the same behaviour for any smooth loss.
\end{remark}

\section{Momentum and adaptive steps}
\labsec{momentum}

Momentum averages successive gradients,
\[
	\vect{v}_{k+1}=\beta\vect{v}_k+\nabla f(\vect{x}_k),
	\qquad
	\vect{x}_{k+1}=\vect{x}_k-\alpha\vect{v}_{k+1},
\]
which cancels the alternating components of the zigzag while accumulating the
consistent ones. Adaptive methods instead rescale each coordinate by a running
measure of its own gradient magnitude, which amounts to approximating the
diagonal of the Hessian on the fly.

\begin{remark}
Both are attempts to reduce the effective condition number without forming
second derivatives. Neither changes the fact, established above, that the
condition number is what has to be reduced.
\end{remark}

\section{Newton's method}
\labsec{newton}

Minimising the quadratic model \refeq{taylor2} exactly, rather than following its
linear part, gives
\[
	\vect{x}_{k+1}=\vect{x}_k-H(\vect{x}_k)^{-1}\nabla f(\vect{x}_k).
\]
On a quadratic objective this reaches the minimum in a single step, and near a
non-degenerate minimum it converges quadratically. The obstacle is cost: forming
and factorising \(H\) is \(O(n^{3})\) work in \(n\) parameters, which is
unavailable when \(n\) is in the millions. Quasi-Newton methods maintain a
cheap approximation instead.

\section{Stochastic gradients}
\labsec{sgd}

Training losses are sums over examples,
\(f(\vect{x})=\frac1n\sum_{i=1}^{n}\ell_i(\vect{x})\). Evaluating the full
gradient costs one pass over the data; evaluating it on a random mini-batch
costs a fraction of that and gives an unbiased estimate, since the expectation of
the batch gradient is the full gradient.

\begin{remark}
The trade is variance for speed. Because the estimate is noisy, a constant step
size leaves the iterates bouncing near the optimum rather than settling, which is
why step sizes are decayed. The unbiasedness claim itself belongs to
\refch{expectation}.
\end{remark}

\section{Backpropagation}
\labsec{backprop}

For a composition \(f=f_L\circ\cdots\circ f_1\) with a scalar output, the chain
rule of \refsec{matrix-chain-rule} gives
\[
	\nabla f=J_1\tp J_2\tp\cdots J_L\tp\,,
\]
a product that may be evaluated in either order. Multiplying from the left
(inputs first) propagates full Jacobians, each of size (layer width) squared.
Multiplying from the right (output first) propagates a single vector at every
stage, because the output is scalar.

Backpropagation is that second ordering, nothing more: a forward pass storing
the intermediate values, then a backward pass carrying one vector through the
transposed Jacobians. Its cost is a small constant multiple of the forward pass,
independent of the number of parameters, and that asymptotic fact is what makes
training large models possible at all.

\marginnote{The name suggests a special algorithm. It is better understood as a
choice of association order in a matrix product --- the choice that keeps every
intermediate object as small as the output rather than as large as the layer.}
